% ============================================================
%  references.tex — consolidated bibliography [1]-[91]
%  Plain thebibliography: reproduces the numbering exactly as
%  proofread in the drafted chapters. Cited via \cite{refN}.
%
%  UPGRADE PATH: to switch to biblatex+biber later (recommended if
%  you keep adding references), convert each \bibitem below into a
%  proper .bib entry — see the chat message for why this was not
%  done automatically.
% ============================================================
\begin{thebibliography}{99}

\bibitem{ref1} H. Arksey and L. O'Malley, "Scoping studies: Towards a methodological framework," *Int. J. Soc. Res. Methodol.*, vol. 8, no. 1, pp. 19–32, Feb. 2005, doi: \url{https://doi.org/10}.1080/1364557032000119616.

\bibitem{ref2} D. O. Kazmer, *Injection Mold Design Engineering*, 2nd ed. Munich, Germany: Hanser, 2016.

\bibitem{ref3} D. V. Rosato and M. G. Rosato, *Injection Molding Handbook*, 3rd ed. Boston, MA, USA: Kluwer Academic, 2000.

\bibitem{ref4} P. Zhao, H. Zhou, Y. Li, and D. Li, "Intelligent injection molding on sensing, optimization, and control," *Adv. Polym. Technol.*, vol. 2020, art. 7023616, pp. 1–22, 2020, doi: \url{https://doi.org/10}.1155/2020/7023616.

\bibitem{ref5} A. Polenta, S. Tomassini, N. Falcionelli, P. Contardo, A. F. Dragoni, and P. Sernani, "A comparison of machine learning techniques for the quality classification of molded products," *Information*, vol. 13, no. 6, art. 272, Jun. 2022, doi: \url{https://doi.org/10}.3390/info13060272.

\bibitem{ref6} *Road Lighting — Part 2: Performance Requirements*, UNI EN 13201-2:2016, European Committee for Standardization, Brussels, Belgium, 2016.

\bibitem{ref7} F. K. Mukarker, "Product and process innovations in the Palestinian plastic industry," M.S. thesis, Faculty of Graduate Studies, Al-Quds University, Jerusalem, Palestine, 2016.

\bibitem{ref8} *Quality Management Systems — Requirements*, ISO 9001:2015, International Organization for Standardization, Geneva, Switzerland, 2015.

\bibitem{ref9} B. Ribeiro, "Support vector machines for quality monitoring in a plastic injection molding process," *IEEE Trans. Syst., Man, Cybern. C, Appl. Rev.*, vol. 35, no. 3, pp. 401–410, Aug. 2005, doi: \url{https://doi.org/10}.1109/TSMCC.2004.843228.

\bibitem{ref10} R. D. Párizs, D. Török, T. Ageyeva, and J. G. Kovács, "Machine learning in injection molding: An Industry 4.0 method of quality prediction," *Sensors*, vol. 22, no. 7, art. 2704, Apr. 2022, doi: \url{https://doi.org/10}.3390/s22072704.

\bibitem{ref11} O. Ogorodnyk, O. V. Lyngstad, M. Larsen, K. Wang, and K. Martinsen, "Application of machine learning methods for prediction of parts quality in thermoplastics injection molding," in *Advanced Manufacturing and Automation VIII*, K. Wang, Y. Wang, J. O. Strandhagen, and T. Yu, Eds. Singapore: Springer, 2019, pp. 237–244, doi: \url{https://doi.org/10}.1007/978-981-13-2375-1\_30.

\bibitem{ref12} B. Silva, R. Marques, D. Faustino, P. Ilheu, T. Santos, J. Sousa, and A. D. Rocha, "Enhance the injection molding quality prediction with artificial intelligence to reach zero-defect manufacturing," *Processes*, vol. 11, no. 1, art. 62, Dec. 2022, doi: \url{https://doi.org/10}.3390/pr11010062.

\bibitem{ref13} Z. Hu, Z. Yin, L. Qin, and F. Xu, "A novel method of fault diagnosis for injection molding systems based on improved VGG16 and machine vision," *Sustainability*, vol. 14, no. 21, art. 14280, Nov. 2022, doi: \url{https://doi.org/10}.3390/su142114280.

\bibitem{ref14} R. X. Gao, X. Tang, G. Gordon, and D. O. Kazmer, "Online product quality monitoring through in-process measurement," *CIRP Ann. Manuf. Technol.*, vol. 63, no. 1, pp. 493–496, 2014, doi: \url{https://doi.org/10}.1016/j.cirp.2014.03.041.

\bibitem{ref15} G. Gordon, D. O. Kazmer, X. Tang, Z. Fan, and R. X. Gao, "Quality control using a multivariate injection molding sensor," *Int. J. Adv. Manuf. Technol.*, vol. 78, no. 9–12, pp. 1381–1391, 2015, doi: \url{https://doi.org/10}.1007/s00170-014-6706-6.

\bibitem{ref16} Q. Wang, X. Zhao, J. Zhang, P. Zhang, X. Wang, C. Yang, J. Wang, and Z. Wu, "Research on quality characterization method of micro-injection products based on cavity pressure," *Polymers*, vol. 13, no. 16, art. 2755, Aug. 2021, doi: \url{https://doi.org/10}.3390/polym13162755.

\bibitem{ref17} S. Kumar, H. S. Park, and C. M. Lee, "Data-driven smart control of injection molding process," *CIRP J. Manuf. Sci. Technol.*, vol. 31, pp. 439–449, 2020, doi: \url{https://doi.org/10}.1016/j.cirpj.2020.07.006.

\bibitem{ref18} H. Tercan, P. Deibert, and T. Meisen, "Continual learning of neural networks for quality prediction in production using memory aware synapses and weight transfer," *J. Intell. Manuf.*, vol. 33, no. 1, pp. 283–292, Jan. 2022, doi: \url{https://doi.org/10}.1007/s10845-021-01793-0.

\bibitem{ref19} M. Asadi, S. S. Hashemi, and M. T. Sadeghi, "Injection molding inspection system based on machine vision," in *Proc. 9th RSI Int. Conf. Robotics and Mechatronics (ICRoM)*, 2021, pp. 536–541, doi: \url{https://doi.org/10}.1109/ICRoM54204.2021.9663455.

\bibitem{ref20} M. P. Muresan, D. G. Cireap, and I. Giosan, "Automatic vision inspection solution for the manufacturing process of automotive components through plastic injection molding," in *Proc. IEEE 16th Int. Conf. Intelligent Computer Communication and Processing (ICCP)*, Sep. 2020, pp. 423–430, doi: \url{https://doi.org/10}.1109/ICCP51029.2020.9266249.

\bibitem{ref21} R. Zhang, Y. Li, and Z. Guan, "Research on injection molded parts defect detection algorithm based on multiplicative feature fusion and improved attention mechanism," *Sci. Rep.*, vol. 14, art. 30894, Dec. 2024, doi: \url{https://doi.org/10}.1038/s41598-024-81430-x.

\bibitem{ref22} Y. Chen, Y. Ding, F. Zhao, E. Zhang, Z. Wu, and L. Shao, "Surface defect detection methods for industrial products: A review," *Appl. Sci.*, vol. 11, no. 16, art. 7657, Aug. 2021, doi: \url{https://doi.org/10}.3390/app11167657.

\bibitem{ref23} S. Kabir, M. S. Hossain, and K. Andersson, "A review of explainable artificial intelligence from the perspectives of challenges and opportunities," *Algorithms*, vol. 18, no. 9, art. 556, Sep. 2025, doi: \url{https://doi.org/10}.3390/a18090556.

\bibitem{ref24} C. Rudin, "Stop explaining black box machine learning models for high stakes decisions and use interpretable models instead," *Nat. Mach. Intell.*, vol. 1, no. 5, pp. 206–215, May 2019, doi: \url{https://doi.org/10}.1038/s42256-019-0048-x.

\bibitem{ref25} S. M. Lundberg and S.-I. Lee, "A unified approach to interpreting model predictions," in *Advances in Neural Information Processing Systems (NeurIPS)*, vol. 30, 2017, pp. 4765–4774.

\bibitem{ref26} S. M. Lundberg, G. Erion, H. Chen, A. DeGrave, J. M. Prutkin, B. Nair, R. Katz, J. Himmelfarb, N. Bansal, and S.-I. Lee, "From local explanations to global understanding with explainable AI for trees," *Nat. Mach. Intell.*, vol. 2, no. 1, pp. 56–67, Jan. 2020, doi: \url{https://doi.org/10}.1038/s42256-019-0138-9.

\bibitem{ref27} M. T. Ribeiro, S. Singh, and C. Guestrin, "'Why should I trust you?': Explaining the predictions of any classifier," in *Proc. 22nd ACM SIGKDD Int. Conf. Knowledge Discovery and Data Mining*, 2016, pp. 1135–1144, doi: \url{https://doi.org/10}.1145/2939672.2939778.

\bibitem{ref28} S. Wachter, B. Mittelstadt, and C. Russell, "Counterfactual explanations without opening the black box: Automated decisions and the GDPR," *Harvard J. Law Technol.*, vol. 31, no. 2, pp. 841–887, 2018.

\bibitem{ref29} R. K. Mothilal, A. Sharma, and C. Tan, "Explaining machine learning classifiers through diverse counterfactual explanations," in *Proc. 2020 Conf. Fairness, Accountability, and Transparency (FAT\* '20)*, 2020, pp. 607–617, doi: \url{https://doi.org/10}.1145/3351095.3372850.

\bibitem{ref30} S. Verma, V. Boonsanong, M. Hoang, K. Hines, J. Dickerson, and C. Shah, "Counterfactual explanations and algorithmic recourses for machine learning: A review," *ACM Comput. Surv.*, vol. 56, no. 12, art. 312, pp. 1–42, Dec. 2024, doi: \url{https://doi.org/10}.1145/3677119.

\bibitem{ref31} R. Guidotti, "Counterfactual explanations and how to find them: Literature review and benchmarking," *Data Min. Knowl. Discov.*, vol. 38, no. 5, pp. 2770–2824, 2024, doi: \url{https://doi.org/10}.1007/s10618-022-00831-6.

\bibitem{ref32} R. Poyiadzi, K. Sokol, R. Santos-Rodriguez, T. De Bie, and P. Flach, "FACE: Feasible and actionable counterfactual explanations," in *Proc. AAAI/ACM Conf. AI, Ethics, and Society (AIES)*, 2020, pp. 344–350, doi: \url{https://doi.org/10}.1145/3375627.3375850.

\bibitem{ref33} A. Van Looveren and J. Klaise, "Interpretable counterfactual explanations guided by prototypes," in *Machine Learning and Knowledge Discovery in Databases (ECML PKDD)*, Cham, Switzerland: Springer, 2021, pp. 650–665, doi: \url{https://doi.org/10}.1007/978-3-030-86520-7\_40.

\bibitem{ref34} B. Ustun, A. Spangher, and Y. Liu, "Actionable recourse in linear classification," in *Proc. Conf. Fairness, Accountability, and Transparency (FAT\* '19)*, 2019, pp. 10–19, doi: \url{https://doi.org/10}.1145/3287560.3287566.

\bibitem{ref35} A.-H. Karimi, B. Schölkopf, and I. Valera, "Algorithmic recourse: From counterfactual explanations to interventions," in *Proc. ACM Conf. Fairness, Accountability, and Transparency (FAccT)*, 2021, pp. 353–362, doi: \url{https://doi.org/10}.1145/3442188.3445899.

\bibitem{ref36} M. T. Keane, E. M. Kenny, E. Delaney, and B. Smyth, "If only we had better counterfactual explanations: Five key deficits to rectify in the evaluation of counterfactual XAI techniques," in *Proc. 30th Int. Joint Conf. Artificial Intelligence (IJCAI-21)*, Survey Track, 2021, pp. 4466–4474, doi: \url{https://doi.org/10}.24963/ijcai.2021/609.

\bibitem{ref37} D. Müller, M. März, S. Scheele, and U. Schmid, "An interactive explanatory AI system for industrial quality control," in *Proc. AAAI Conf. Artificial Intelligence*, vol. 36, no. 11, 2022, pp. 12580–12586, doi: \url{https://doi.org/10}.1609/aaai.v36i11.21530.

\bibitem{ref38} R. Hoffmann and C. Reich, "A systematic literature review on artificial intelligence and explainable artificial intelligence for visual quality assurance in manufacturing," *Electronics*, vol. 12, no. 22, art. 4572, Nov. 2023, doi: \url{https://doi.org/10}.3390/electronics12224572.

\bibitem{ref39} A. Saadallah, J. Büscher, O. Abdulaaty, T. Panusch, J. Deuse, and K. Morik, "Explainable predictive quality inspection using deep learning in electronics manufacturing," *Procedia CIRP*, vol. 107, pp. 594–599, 2022, doi: \url{https://doi.org/10}.1016/j.procir.2022.05.031.

\bibitem{ref40} A. Ettalibi, A. Elouadi, and A. Mansour, "AI and computer vision-based real-time quality control: A review of industrial applications," *Procedia Comput. Sci.*, vol. 231, pp. 212–220, 2024, doi: \url{https://doi.org/10}.1016/j.procs.2023.12.195.

\bibitem{ref41} M. Carletti, C. Masiero, A. Beghi, and G. A. Susto, "Explainable machine learning in Industry 4.0: Evaluating feature importance in anomaly detection to enable root cause analysis," in *Proc. IEEE Int. Conf. Systems, Man and Cybernetics (SMC)*, 2019, pp. 21–26, doi: \url{https://doi.org/10}.1109/SMC.2019.8913901.

\bibitem{ref42} S. K. Selvaraj, A. Raj, R. R. Mahadevan, U. Chadha, and V. Paramasivam, "A review on machine learning models in injection molding machines," *Adv. Mater. Sci. Eng.*, vol. 2022, art. 1949061, 2022, doi: \url{https://doi.org/10}.1155/2022/1949061.

\bibitem{ref43} H. Lee, K. Ryu, and Y. Cho, "A framework of a smart injection molding system based on real-time data," *Procedia Manuf.*, vol. 11, pp. 1004–1011, 2017, doi: \url{https://doi.org/10}.1016/j.promfg.2017.07.206.

\bibitem{ref44} J. Gim and B. Rhee, "Novel analysis methodology of cavity pressure profiles in injection-molding processes using interpretation of machine learning model," *Polymers*, vol. 13, no. 19, art. 3297, Oct. 2021, doi: \url{https://doi.org/10}.3390/polym13193297.

\bibitem{ref45} J. Gim and L.-S. Turng, "Interpretation of the effect of transient process data on part quality of injection molding based on explainable artificial intelligence," *Int. J. Prod. Res.*, vol. 61, no. 23, pp. 8192–8212, 2023, doi: \url{https://doi.org/10}.1080/00207543.2022.2153185.

\bibitem{ref46} J. Gim, C.-Y. Lin, and L.-S. Turng, "In-mold condition-centered and explainable artificial intelligence-based (IMC-XAI) process optimization for injection molding," *J. Manuf. Syst.*, vol. 72, pp. 196–213, Feb. 2024, doi: \url{https://doi.org/10}.1016/j.jmsy.2023.11.017.

\bibitem{ref47} M. Muaz, S. Sajid, T. Schulze, C. Liu, N. Klasen, and B. Drescher, "Explainable AI for correct root cause analysis of product quality in injection moulding," *J. Manuf. Process.*, vol. 145, pp. 371–380, Jul. 2025, doi: \url{https://doi.org/10}.1016/j.jmapro.2025.03.114.

\bibitem{ref48} G. Marín Díaz, "Comparative analysis of explainable AI methods for manufacturing defect prediction: A mathematical perspective," *Mathematics*, vol. 13, no. 15, art. 2436, Aug. 2025, doi: \url{https://doi.org/10}.3390/math13152436.

\bibitem{ref49} G. Brue, *Six Sigma for Managers*, 1st ed. New York, NY, USA: McGraw-Hill, 2002.

\bibitem{ref50} K. Papageorgiou, T. Theodosiou, A. Rapti, E. I. Papageorgiou, N. Dimitriou, D. Tzovaras, and G. Margetis, "A systematic review on machine learning methods for root cause analysis towards zero-defect manufacturing," *Front. Manuf. Technol.*, vol. 2, art. 972712, Oct. 2022, doi: \url{https://doi.org/10}.3389/fmtec.2022.972712.

\bibitem{ref51} M. Solé, V. Muntés-Mulero, A. I. Rana, and G. Estrada, "Survey on models and techniques for root-cause analysis," arXiv:1701.08546, Jul. 2017.

\bibitem{ref52} E. e Oliveira, V. L. Miguéis, and J. L. Borges, "Automatic root cause analysis in manufacturing: An overview and conceptualization," *J. Intell. Manuf.*, vol. 34, no. 5, pp. 2061–2078, 2023, doi: \url{https://doi.org/10}.1007/s10845-022-01914-3.

\bibitem{ref53} A. Chigurupati and N. Lassar, "Root cause analysis using artificial intelligence," in *Proc. Annu. Reliability and Maintainability Symp. (RAMS)*, 2017, pp. 1–5, doi: \url{https://doi.org/10}.1109/RAM.2017.7889651.

\bibitem{ref54} A. Lokrantz, E. Gustavsson, and M. Jirstrand, "Root cause analysis of failures and quality deviations in manufacturing using machine learning," *Procedia CIRP*, vol. 72, pp. 1057–1062, 2018, doi: \url{https://doi.org/10}.1016/j.procir.2018.03.229.

\bibitem{ref55} K. Wang, C. Liu, and Y. Lu, "Ensemble Bayesian network for root cause analysis of product defects via learning from historical production data," *J. Manuf. Syst.*, vol. 75, pp. 102–115, Aug. 2024, doi: \url{https://doi.org/10}.1016/j.jmsy.2024.06.001.

\bibitem{ref56} C. Wehner, M. Kertel, and J. Wewerka, "Interactive and intelligent root cause analysis in manufacturing with causal Bayesian networks and knowledge graphs," in *Proc. IEEE 97th Vehicular Technology Conf. (VTC2023-Spring)*, 2023, pp. 1–7, doi: \url{https://doi.org/10}.1109/VTC2023-Spring57618.2023.10199975.

\bibitem{ref57} Q. Ma, H. Li, and A. Thorstenson, "A big data-driven root cause analysis system: Application of machine learning in quality problem solving," *Comput. Ind. Eng.*, vol. 160, art. 107580, Oct. 2021, doi: \url{https://doi.org/10}.1016/j.cie.2021.107580.

\bibitem{ref58} D. J. Huang and H. Li, "A machine learning guided investigation of quality repeatability in metal laser powder bed fusion additive manufacturing," *Mater. Des.*, vol. 203, art. 109606, May 2021, doi: \url{https://doi.org/10}.1016/j.matdes.2021.109606.

\bibitem{ref59} S. Sarkar, N. Ejaz, M. Kumar, and J. Maiti, "Root cause analysis of incidents using text clustering and classification algorithms," in *Proc. ICETIT 2019*, Lecture Notes in Electrical Engineering, vol. 605. Cham, Switzerland: Springer, 2020, pp. 707–718, doi: \url{https://doi.org/10}.1007/978-3-030-30577-2\_63.

\bibitem{ref60} J. D. Crocco and P. R. O'Hern, Jr., "Manufacturing quality improvement through statistical root cause analysis using convolution neural networks," U.S. Patent Application US 2018/0293722 A1, Oct. 11, 2018.

\bibitem{ref61} B. Steenwinckel, D. De Paepe, S. Vanden Hautte, P. Heyvaert, M. Bentefrit, P. Moens, A. Dimou, B. Van Den Bossche, F. De Turck, S. Van Hoecke, and F. Ongenae, "FLAGS: A methodology for adaptive anomaly detection and root cause analysis on sensor data streams by fusing expert knowledge with machine learning," *Future Gener. Comput. Syst.*, vol. 116, pp. 30–48, Mar. 2021, doi: \url{https://doi.org/10}.1016/j.future.2020.10.015.

\bibitem{ref62} T. Laugel, M.-J. Lesot, C. Marsala, X. Renard, and M. Detyniecki, "The dangers of post-hoc interpretability: Unjustified counterfactual explanations," in *Proc. 28th Int. Joint Conf. Artificial Intelligence (IJCAI-19)*, 2019, pp. 2801–2807, doi: \url{https://doi.org/10}.24963/ijcai.2019/388.

\bibitem{ref63} C. A. Escobar and R. Morales-Menendez, *Machine Learning in Manufacturing: Quality 4.0 and the Zero Defects Vision*, 1st ed. Amsterdam, The Netherlands: Elsevier, 2024.

\bibitem{ref64} C. A. Escobar, M. E. McGovern, and R. Morales-Menendez, "Quality 4.0: A review of big data challenges in manufacturing," *J. Intell. Manuf.*, vol. 32, no. 8, pp. 2319–2334, 2021, doi: \url{https://doi.org/10}.1007/s10845-021-01765-4.

\bibitem{ref65} P. Bibow, M. Dalibor, C. Hopmann, B. Mainz, B. Rumpe, D. Schmalzing, M. Schmitz, and A. Wortmann, "Model-driven development of a digital twin for injection molding," in *Advanced Information Systems Engineering (CAiSE 2020)*, Lecture Notes in Computer Science, vol. 12127. Cham, Switzerland: Springer, 2020, pp. 85–100, doi: \url{https://doi.org/10}.1007/978-3-030-49435-3\_6.

\bibitem{ref66} D. C. Montgomery, *Introduction to Statistical Quality Control*, 8th ed. Hoboken, NJ, USA: Wiley, 2019.

\bibitem{ref67} J. Gama, I. Žliobaitė, A. Bifet, M. Pechenizkiy, and A. Bouchachia, "A survey on concept drift adaptation," *ACM Comput. Surv.*, vol. 46, no. 4, art. 44, pp. 1–37, Mar. 2014, doi: \url{https://doi.org/10}.1145/2523813.

\bibitem{ref68} J. Lu, A. Liu, F. Dong, F. Gu, J. Gama, and G. Zhang, "Learning under concept drift: A review," *IEEE Trans. Knowl. Data Eng.*, vol. 31, no. 12, pp. 2346–2363, Dec. 2019, doi: \url{https://doi.org/10}.1109/TKDE.2018.2876857.

\bibitem{ref69} B. Yurdakul, "Statistical properties of population stability index," Ph.D. dissertation, Western Michigan Univ., Kalamazoo, MI, USA, 2018.

\bibitem{ref70} M. Nauta, J. Trienes, S. Pathak, E. Nguyen, M. Peters, Y. Schmitt, J. Schlötterer, M. van Keulen, and C. Seifert, "From anecdotal evidence to quantitative evaluation methods: A systematic review on evaluating explainable AI," *ACM Comput. Surv.*, vol. 55, no. 13s, art. 295, pp. 1–42, Jul. 2023, doi: \url{https://doi.org/10}.1145/3583558.

\bibitem{ref71} D. Alvarez-Melis and T. S. Jaakkola, "On the robustness of interpretability methods," arXiv:1806.08049, Jun. 2018.

\bibitem{ref72} D. Slack, S. Hilgard, E. Jia, S. Singh, and H. Lakkaraju, "Fooling LIME and SHAP: Adversarial attacks on post hoc explanation methods," in *Proc. AAAI/ACM Conf. AI, Ethics, and Society (AIES)*, 2020, pp. 180–186, doi: \url{https://doi.org/10}.1145/3375627.3375830.

\bibitem{ref73} S. Krishna, T. Han, A. Gu, J. Pombra, S. Jabbari, S. Wu, and H. Lakkaraju, "The disagreement problem in explainable machine learning: A practitioner's perspective," arXiv:2202.01602, Feb. 2022.

\bibitem{ref74} J. Adebayo, J. Gilmer, M. Muelly, I. Goodfellow, M. Hardt, and B. Kim, "Sanity checks for saliency maps," in *Advances in Neural Information Processing Systems (NeurIPS)*, vol. 31, 2018, pp. 9505–9515.

\bibitem{ref75} F. Doshi-Velez and B. Kim, "Towards a rigorous science of interpretable machine learning," arXiv:1702.08608, Mar. 2017.

\bibitem{ref76} M. Pawelczyk, S. Bielawski, J. van den Heuvel, T. Richter, and G. Kasneci, "CARLA: A Python library to benchmark algorithmic recourse and counterfactual explanation algorithms," arXiv:2108.00783, Aug. 2021.

\bibitem{ref77} A. R. Hevner, S. T. March, J. Park, and S. Ram, "Design science in information systems research," *MIS Quarterly*, vol. 28, no. 1, pp. 75–105, Mar. 2004, doi: \url{https://doi.org/10}.2307/25148625.

\bibitem{ref78} K. Peffers, T. Tuunanen, M. A. Rothenberger, and S. Chatterjee, "A design science research methodology for information systems research," *J. Manage. Inf. Syst.*, vol. 24, no. 3, pp. 45–77, 2007, doi: \url{https://doi.org/10}.2753/MIS0742-1222240302.

\bibitem{ref79} L. Breiman, "Random forests," *Mach. Learn.*, vol. 45, no. 1, pp. 5–32, Oct. 2001, doi: \url{https://doi.org/10}.1023/A:1010933404324.

\bibitem{ref80} T. Chen and C. Guestrin, "XGBoost: A scalable tree boosting system," in *Proc. 22nd ACM SIGKDD Int. Conf. Knowledge Discovery and Data Mining*, 2016, pp. 785–794, doi: \url{https://doi.org/10}.1145/2939672.2939785.

\bibitem{ref81} F. Pedregosa, G. Varoquaux, A. Gramfort, V. Michel, B. Thirion, O. Grisel, M. Blondel, P. Prettenhofer, R. Weiss, V. Dubourg, J. Vanderplas, A. Passos, D. Cournapeau, M. Brucher, M. Perrot, and É. Duchesnay, "Scikit-learn: Machine learning in Python," *J. Mach. Learn. Res.*, vol. 12, pp. 2825–2830, 2011.

\bibitem{ref82} P. Geurts, D. Ernst, and L. Wehenkel, "Extremely randomized trees," *Mach. Learn.*, vol. 63, no. 1, pp. 3–42, Apr. 2006, doi: \url{https://doi.org/10}.1007/s10994-006-6226-1.

\bibitem{ref83} R. Kohavi, "A study of cross-validation and bootstrap for accuracy estimation and model selection," in *Proc. 14th Int. Joint Conf. Artificial Intelligence (IJCAI)*, vol. 2, 1995, pp. 1137–1145.

\bibitem{ref84} G. C. Cawley and N. L. C. Talbot, "On over-fitting in model selection and subsequent selection bias in performance evaluation," *J. Mach. Learn. Res.*, vol. 11, pp. 2079–2107, 2010.

\bibitem{ref85} E. S. Page, "Continuous inspection schemes," *Biometrika*, vol. 41, no. 1–2, pp. 100–115, Jun. 1954, doi: \url{https://doi.org/10}.1093/biomet/41.1-2.100.

\bibitem{ref86} S. Seabold and J. Perktold, "Statsmodels: Econometric and statistical modeling with Python," in *Proc. 9th Python in Science Conf. (SciPy)*, 2010, pp. 92–96, doi: \url{https://doi.org/10}.25080/Majora-92bf1922-011.

\bibitem{ref87} Q. McNemar, "Note on the sampling error of the difference between correlated proportions or percentages," *Psychometrika*, vol. 12, no. 2, pp. 153–157, Jun. 1947, doi: \url{https://doi.org/10}.1007/BF02295996.

\bibitem{ref88} T. G. Dietterich, "Approximate statistical tests for comparing supervised classification learning algorithms," *Neural Comput.*, vol. 10, no. 7, pp. 1895–1923, Oct. 1998, doi: \url{https://doi.org/10}.1162/089976698300017197.

\bibitem{ref89} R. Parasuraman and V. Riley, "Humans and automation: Use, misuse, disuse, abuse," *Human Factors*, vol. 39, no. 2, pp. 230–253, Jun. 1997, doi: \url{https://doi.org/10}.1518/001872097778543886.

\bibitem{ref90} J. D. Lee and K. A. See, "Trust in automation: Designing for appropriate reliance," *Human Factors*, vol. 46, no. 1, pp. 50–80, 2004, doi: \url{https://doi.org/10}.1518/hfes.46.1.50\_30392.

\bibitem{ref91} Z. Buçinca, M. B. Malaya, and K. Z. Gajos, "To trust or to think: Cognitive forcing functions can reduce overreliance on AI in AI-assisted decision-making," *Proc. ACM Hum.-Comput. Interact.*, vol. 5, no. CSCW1, art. 188, pp. 1–21, Apr. 2021, doi: \url{https://doi.org/10}.1145/3449287.

\end{thebibliography}